\section{The Verified Algorithm, Level by Level}\label{sec:results}

This section follows Fig.~\ref{fig:stack} from the bottom. Each major result is stated informally
first and then as the Isabelle text, copied from the sources and checked against them
mechanically (Appendix~\ref{app:method}).

\subsection{The Foundation}\label{sec:foundation}

Vern Paulsen's course \emph{Lectures on Quantum Computing and Completely Positive
Maps}~\cite{paulsen2011notes}, given in Houston in 2011, runs over 34 lectures from Hilbert spaces to
quantum error correction. Its formalisation follows the lectures in order, in 46 theories, on top of
a session of reusable matrix analysis, and every definition and theorem of the notes was audited
against the formal text. The audit found false statements in the notes, each witnessed formally:
the Day~29 ``Hamming $[7,1,3]$'' code has distance~2 and corrects no error
(\isa{notes\_hamming\_not\_correcting}), and on Day~32 two Pauli errors that the notes call
distinguishable agree on the Shor code space (\isa{shor\_notes\_Z1\_Z2}). For HHL, the first ten
lectures matter most: the spectral theorem, positivity and square roots, tensor products and
measurement. The later lectures give the meaning of quantum programs: Choi's theorem
(Listing~\ref{lst:choi}) characterises completely positive maps, which is what a quantum program
denotes, and a theory of the HHL entry reads off Kraus operators from any loop-free program and proves
the program a channel.

Two bridges connect the lectures to the rest of the AFP. Formalising a lecture course means
defining what the course defines, so the notes have their own Hermitian, positive semidefinite and
density matrices; the theory \isa{Notes\_Library\_Bridge} proves each equal to its counterpart in
QHLProver and \IMD{} (\isa{psd\_AFP}: \isa{psd P $\longleftrightarrow$ positive P}). The session
\isa{Kraus\_Maps\_Bridge} connects the course to Unruh's \emph{Kraus Maps}
entry~\cite{unruh2025kraus}, which develops Kraus maps on trace-class operators of arbitrary Hilbert
spaces but does not prove Choi's theorem. For finite types the two settings meet: a list of Kraus
matrices becomes a Kraus family whose map is the lectures' $\rho\mapsto\sum_iW_i\rho W_i^\dagger$, and
with the lectures' Choi theorem every completely positive map in Paulsen's sense is a Kraus map of the
library (Fig.~\ref{fig:krausbridge}, Listing~\ref{lst:choi}).

\begin{figure}[t]
\centering
\begin{tikzcd}[column sep=12.5em, row sep=3.6em, cells={font=\small},
  every label/.append style={font=\footnotesize}]
\rho\in\C^{n\times n} \arrow[r, "\rho\;\mapsto\;\textstyle\sum_iW_i\rho W_i^\dagger"] \arrow[d, "\mathit{tc\_of\_mat}"']
  & \C^{d\times d} \\
T_\rho\in\mathcal T(\ell^2(\alpha)) \arrow[r, "\mathit{kf\_apply}\;(\mathit{kf\_of\_mats}\;\mathit{Ws})"', shorten=4pt]
  & \mathcal T(\ell^2(\beta)) \arrow[u, "\mathit{mat\_of\_tc}"']
\end{tikzcd}
\caption{The bridge to \emph{Kraus Maps} as a commuting square (\isa{kf\_of\_mats\_apply}): for finite types
$\alpha$, $\beta$ of sizes $n$ and $d$, the lectures' Kraus map on matrices (top) agrees with Unruh's Kraus map
on trace-class operators (bottom), read through the translation between matrices and operators. With
Choi's theorem it follows that every completely positive map of the lectures is a Kraus map of the
library (\isa{cp\_kraus\_map}).}\label{fig:krausbridge}
\end{figure}

%%LST Day19_Choi_Theorem.thy choi_theorem ; Kraus_Maps_Bridge/Kraus_Maps_Bridge.thy cp_kraus_map
\begin{lstlisting}[float=tp,caption={Choi's theorem (Lecture 19), and the lectures' completely positive maps as Kraus maps of the AFP entry \emph{Kraus Maps}; \isa{tc\_map} reads a map of matrices as a map of trace-class operators on \isa{'a ell2}},label={lst:choi}]
theorem choi_theorem:
  fixes ~$\Phi$~ :: "complex Matrix.mat ~$\Rightarrow$~ complex Matrix.mat"
  assumes linear: "~$\bigwedge$~A B. A ~$\in$~ carrier_mat n n ~$\Longrightarrow$~ B ~$\in$~ carrier_mat n n
      ~$\Longrightarrow$~ ~$\Phi$~ (A + B) = ~$\Phi$~ A + ~$\Phi$~ B"
    and smult: "~$\bigwedge$~c A. A ~$\in$~ carrier_mat n n ~$\Longrightarrow$~ ~$\Phi$~ (c ~$\cdot_{m}$~ A) = c ~$\cdot_{m}$~ ~$\Phi$~ A"
    and range: "~$\forall$~A. A ~$\in$~ carrier_mat n n ~$\longrightarrow$~ ~$\Phi$~ A ~$\in$~ carrier_mat d d"
    and npos: "n > 0"
    and dpos: "d > 0"
  shows "completely_positive n d ~$\Phi$~ ~$\longleftrightarrow$~ psd (choi_matrix n d ~$\Phi$~)"
theorem cp_kraus_map:
  fixes ~$\Phi$~ :: "complex Matrix.mat ~$\Rightarrow$~ complex Matrix.mat"
  assumes linear: "~$\bigwedge$~A B. A ~$\in$~ carrier_mat CARD('a) CARD('a) ~$\Longrightarrow$~ B ~$\in$~
      carrier_mat CARD('a) CARD('a) ~$\Longrightarrow$~ ~$\Phi$~ (A + B) = ~$\Phi$~ A + ~$\Phi$~ B"
    and smult: "~$\bigwedge$~c A. A ~$\in$~ carrier_mat CARD('a) CARD('a) ~$\Longrightarrow$~ ~$\Phi$~ (c ~$\cdot_{m}$~ A)
        = c ~$\cdot_{m}$~ ~$\Phi$~ A"
    and range: "~$\forall$~A. A ~$\in$~ carrier_mat CARD('a) CARD('a) ~$\longrightarrow$~ ~$\Phi$~ A ~$\in$~
        carrier_mat CARD('b) CARD('b)"
    and cp: "completely_positive CARD('a) CARD('b) ~$\Phi$~"
  shows "Kraus_Maps.kraus_map (tc_map ~$\Phi$~ :: ('a::enum ell2, 'a ell2)
      trace_class ~$\Rightarrow$~ ('b::enum ell2, 'b ell2) trace_class)"
\end{lstlisting}

\subsection{Mathematics}\label{sec:math}

\paragraph{A functional calculus.} HHL computes $A^{-1}$, $e^{iAt}$ and $\sqrt{I-C^2A^{-2}}$ of a
Hermitian matrix, so its proof needs $g(A)$ for real functions~$g$. We define $g(A)$ by what it does
to eigenvectors rather than through a chosen diagonalisation.

\begin{theorem}[functional calculus]\label{thm:matfun}
Let $A\in\C^{n\times n}$ be Hermitian. There is exactly one matrix $g(A)$ with $g(A)v=g(\lambda)v$ for
every eigenvector $v$ of $A$ with eigenvalue~$\lambda$, and $g(A)=U\operatorname{diag}(g(d_1),\dots,g(d_n))
U^\dagger$ for every spectral decomposition $U^\dagger AU=\operatorname{diag}(d_1,\dots,d_n)$.\isalink{Hermitian_Functional_Calculus/Hermitian_Functional_Calculus.thy}{497}
\end{theorem}

%%LST Hermitian_Functional_Calculus/Hermitian_Functional_Calculus.thy acts_on_eigenvectors mat_fun_unique
\begin{lstlisting}[float=tp,caption={The functional calculus: \isa{B acts\_as g on\_eigenvectors\_of A} is notation for \isa{acts\_on\_eigenvectors A g B}, and \isa{mat\_fun g A} is the unique such \isa{B}},label={lst:matfun}]
definition acts_on_eigenvectors :: "complex Matrix.mat ~$\Rightarrow$~ (real ~$\Rightarrow$~
    complex) ~$\Rightarrow$~ complex Matrix.mat ~$\Rightarrow$~ bool" where
  "acts_on_eigenvectors A g B ~$\longleftrightarrow$~ (~$\forall$~v l. dim_vec v = dim_row A ~$\longrightarrow$~ A *~$_{v}$~
      v = complex_of_real l ~$\cdot_{v}$~ v ~$\longrightarrow$~ B *~$_{v}$~ v = g l ~$\cdot_{v}$~ v)"
corollary mat_fun_unique:
  assumes A: "A ~$\in$~ carrier_mat n n" "hermitian A"
  shows "~$\exists$~!B. B ~$\in$~ carrier_mat n n ~$\wedge$~ B acts_as g on_eigenvectors_of A"
\end{lstlisting}

The definition makes the algebra of the calculus short: two matrices that agree on all
eigenvectors of~$A$ are equal, because the eigenvectors span, so identities such as
$g(A)h(A)=(gh)(A)$ and $g(A)^\dagger=\bar g(A)$ need no sums over an eigenbasis. The same principle
proves HHL itself (Section~\ref{sec:algorithm}). On top of the calculus we verify the Lie--Trotter
formula with an explicit bound,
\[
\norm{(e^{itH_1/r}\cdots e^{itH_m/r})^r-e^{it(H_1+\dots+H_m)}}\le2(mLt)^2/r
\]
for Hermitian $H_j$ of norm at most~$L$, and Suzuki's higher-order formulas with error of order
$|t|^{2k+3}/r^{2k+2}$.\isalink{Hermitian_Functional_Calculus/Trotter_Product_Formula.thy}{612}

\paragraph{Phase estimation.} Phase estimation is HHL's engine. With $F_N$ the Fourier matrix and
$\Lambda(U)=\sum_k\ketbra kk\otimes U^k$:

\begin{theorem}[phase estimation]\label{thm:qpe}
If $Uv=e^{2\pi i\varphi}v$, then $(F_N^\dagger\otimes I)\Lambda(U)(u\otimes v)=(\sum_j\alpha_j\ket
j)\otimes v$ with $\alpha_j=\frac1N\sum_{k<N}e^{2\pi ik(\varphi-j/N)}$. If $\varphi=j_0/N$, then
$\alpha=\ket{j_0}$. If $|\varphi-j/N|\le1/(2N)$, then $|\alpha_j|^2\ge4/\pi^2$.\isalink{Algorithms/Phase_Estimation.thy}{287}
\end{theorem}

%%LST Algorithms/Phase_Estimation.thy phase_estimation phase_estimation_bound
\begin{lstlisting}[float=tp,caption={Phase estimation and the $4/\pi^2$ bound},label={lst:qpe}]
theorem phase_estimation:
  assumes U: "U ~$\in$~ carrier_mat m m" and u: "dim_vec u = m"
    and ev: "U *~$_{v}$~ u = cis (2 * pi * ~$\varphi$~) ~$\cdot_{v}$~ u"
  shows "(((fourier N)~$^{\dagger}$~) ~$\otimes$~ 1~$_{m}$~ m) *~$_{v}$~ (ctrl_pow N m U *~$_{v}$~ tensor_vec
      (uniform_state N) u) = tensor_vec (Matrix.vec N (pe_amp N ~$\varphi$~)) u"
theorem phase_estimation_bound:
  assumes N: "0 < N" and close: "~$|\varphi$~ - real j / real N~$|$~ ~$\le$~ 1 / (2 * real
      N)"
  shows "4 / pi~$^{2}$~ ~$\le$~ (cmod (pe_amp N ~$\varphi$~ j))~$^{2}$~"
\end{lstlisting}

We also prove that the Fourier circuit of the AFP entry~\cite{manrique2025qft} computes $F_N$, so
the theorem holds for an actual circuit, and the textbook precision bound: with
$n+\lceil\log_2(2+1/(2\varepsilon))\rceil$ clock qubits the estimate is within $2^{-n}$ with probability
at least $1-\varepsilon$~\cite{nielsenchuang}. As a program, phase estimation satisfies, for every set
$J$ of clock values, the triple $\hoaret{(\sum_{j\in J}|\alpha_j|^2)\,I}{\mathit{QPE}}{\Pi_J}$, whose
precondition is the weakest possible.

\begin{figure}[t]
\centering
\begin{tikzpicture}
\begin{axis}[width=11.2cm, height=4.6cm, ymode=log, log origin y=infty, ymin=1e-7, ymax=1.5,
  xmin=-0.6, xmax=15.6, axis lines=left, xtick={0,...,15}, ytick={1e-6,1e-4,1e-2,1},
  xlabel={clock outcome $j$}, ylabel={$|\alpha_j|^2$}, tick label style={font=\scriptsize},
  label style={font=\footnotesize}, legend style={font=\scriptsize, draw=none, fill=white, fill opacity=0.85,
  text opacity=1, at={(0.99,0.99)}, anchor=north east}, legend cell align=left]
\fill[red!7] (axis cs:-0.6,1e-7) rectangle (axis cs:2.3,1.5);
\fill[red!7] (axis cs:8.3,1e-7) rectangle (axis cs:15.6,1.5);
\addplot[ycomb, very thick, hhl!80, mark=*, mark size=1.1pt] coordinates
  {(0,0.003435) (1,0.004576) (2,0.007018) (3,0.01342) (4,0.0401) (5,0.7377) (6,0.1362) (7,0.02382)
   (8,0.009999) (9,0.005794) (10,0.004022) (11,0.003158) (12,0.002731) (13,0.002566) (14,0.002606) (15,0.002864)};
\addlegendentry{uniform clock}
\addplot[only marks, hfc!90!black, mark=square*, mark size=1.5pt] coordinates
  {(0,1.228e-05) (1,3.794e-05) (2,0.0001311) (3,0.0006419) (4,0.008282) (5,0.6851) (6,0.3029) (7,0.002423)
   (8,0.0003187) (9,7.778e-05) (10,2.41e-05) (11,7.718e-06) (12,1.925e-06) (13,9.227e-08) (14,5.157e-07) (15,3.57e-06)};
\addlegendentry{sine-window clock}
\addplot[densely dotted, thick, hhl!80, domain=-0.6:4.3, samples=40, forget plot] {1/(4*(5.3-x)^2)};
\addplot[densely dotted, thick, hhl!80, domain=6.3:13.3, samples=40, forget plot] {1/(4*(x-5.3)^2)};
\addplot[densely dotted, thick, hhl!80, domain=13.3:15.6, samples=10, forget plot] {1/(4*(21.3-x)^2)};
\addplot[densely dashed, thick, hfc!90!black, domain=-0.6:4.3, samples=40, forget plot] {pi^2/(2*(5.3-x)^4)};
\addplot[densely dashed, thick, hfc!90!black, domain=6.3:13.3, samples=40, forget plot] {pi^2/(2*(x-5.3)^4)};
\addplot[densely dashed, thick, hfc!90!black, domain=13.3:15.6, samples=10, forget plot] {pi^2/(2*(21.3-x)^4)};
\addplot[gray, domain=-0.6:15.6, samples=2, forget plot] {4/pi^2};
\node[font=\scriptsize, gray, anchor=south west] at (axis cs:9.2,0.42) {$4/\pi^2$};
\end{axis}
\end{tikzpicture}
\caption{The outcome distribution of phase estimation for $N=16$ and $N\varphi=5.3$ (logarithmic scale),
computed from the definitions, with the two clock states. The nearest outcome has probability at least
$4/\pi^2$ (Theorem~\ref{thm:qpe}). Away from it, the uniform clock's amplitudes obey only
$|\alpha_j|^2\le1/(4N^2\delta_j^2)$ (dotted, \isa{pe\_amp\_tail}), the window's obey
$|\alpha_j|^2\le\pi^2/(2N^4\delta_j^4)$ (dashed, \isa{win\_amp\_tail}). The outcomes more than three
steps from $N\varphi$ (shaded) have total probability $0.039$ with the uniform clock and $0.0003$ with
the window: this is why the window makes $N\sim1/\varepsilon$ enough (Theorem~\ref{thm:window}).}\label{fig:qpe}
\end{figure}

\subsection{The Algorithm}\label{sec:algorithm}

\paragraph{Exact HHL.} Let $W=(F^\dagger\otimes I)\,\Lambda(U)$ be phase estimation without its
input, and $R$ the clock-controlled rotation.

\begin{theorem}[exact HHL]\label{thm:hhl}
Let $A\in\C^{n\times n}$ be Hermitian with every eigenvalue a grid value $\lambda_k$, $0<k<N$. Then
for every $b\in\C^n$ and every real $C$,
\[
W^\dagger RW\,(u\otimes b\otimes\ket0)=u\otimes\bigl(\sqrt{I-C^2A^{-2}}\,b\otimes\ket0+CA^{-1}b\otimes\ket1\bigr).
\]
Hence the ancilla reads $1$ with probability $C^2\norm{A^{-1}b}^2$, and if $C\le|\lambda|\le\Lambda$ on
the spectrum, this is at least $(C/\Lambda)^2\norm b^2$, which is $1/\kappa^2$ for unit
$b$ and $C=\lambda_{\min}$.\isalink{Algorithms/HHL.thy}{590}
\end{theorem}

%%LST Algorithms/HHL.thy hhl
\begin{lstlisting}[float=tp,caption={Exact HHL as a matrix identity, for every right-hand side},label={lst:hhl}]
theorem hhl:
  assumes A: "A ~$\in$~ carrier_mat n n" "hermitian A" and t0: "0 < t0" and
      N: "0 < N"
    and spec: "~$\bigwedge$~l. l ~$\in$~ herm_spectrum A ~$\Longrightarrow$~ ~$\exists$~k. 0 < k ~$\wedge$~ k < N ~$\wedge$~ l =
        reg_eig t0 N k"
    and b: "dim_vec b = n"
  shows "hhl_circuit A t0 N C *~$_{v}$~ tensor_vec (uniform_state N)
      (tensor_vec b (unit_vec 2 0)) = tensor_vec (uniform_state N)
      (tensor_vec (mat_fun (~$\lambda$~x. complex_of_real (sqrt (1 - (C / x)~$^{2}$~)))
      A *~$_{v}$~ b) (unit_vec 2 0) + tensor_vec ((complex_of_real C ~$\cdot_{m}$~
      herm_inv A) *~$_{v}$~ b) (unit_vec 2 1))"
\end{lstlisting}

The proof is the eigenvector principle again. On an eigenvector $v$ of $A$ with eigenvalue
$\lambda_k$, phase estimation writes $\ket k$ exactly, the rotation multiplies by
$R(C/\lambda_k)$, and uncomputation restores the clock; both sides of the identity are linear in $b$
and agree on eigenvectors, so they agree everywhere. No eigenbasis expansion of $b$ appears anywhere
in the proof, although every textbook account is written as one.

\paragraph{Inexact HHL.} The paper's error analysis drops the grid assumption. If $\lambda$ is not a
grid value, phase estimation spreads each eigenvector over the clock, the rotation acts on every
branch with its own angle, and uncomputation is inexact. Branches are orthogonal, which gives an exact
identity for the error on an eigenvector, and the errors of distinct eigenvectors are orthogonal too, so
the bound extends to every $b$ without a dependence on the dimension. Everything then depends on how fast
the clock amplitudes decay away from the true phase. With the uniform clock state they decay like
$1/(N\delta)$, which would need $N\sim1/\varepsilon^2$; the paper prepares the clock in a sine window
instead, whose amplitudes decay like $1/(N\delta)^2$ (Fig.~\ref{fig:qpe}).

\begin{theorem}[window register]\label{thm:window}
Let $0<C$ and $0<\theta<1$, and let $A$ be Hermitian with every eigenvalue in the
interval from $\lambda_{\min}\ge2C/\theta$ to $\pi/(2t_0)$. For every $b$, the circuit with the sine-window clock and the rotation clamped
to $[-1,1]$ satisfies
\[
\norm{\mathrm{out}-\psi\otimes w}^2\le\Bigl(\frac{K_1}{N^2}+\frac{K_2}{N^3}\Bigr)\norm b^2,
\]
with $K_1=8\pi^2(2+\pi^2)\theta^4/(t_0^2C^2(1-\theta^2))$ and $K_2=512\pi^6/(t_0^4\lambda_{\min}^4)$,
where $\psi$ is the window state and $w$ the exact output of Theorem~\ref{thm:hhl}. Hence a clock of size
$N\sim1/\varepsilon$ gives error at most $\varepsilon\norm b$.\isalink{Algorithms/HHL_Window.thy}{739}
\end{theorem}

The paper's own analysis also flags ill-conditioned eigenvalues instead of inverting them. Its
\emph{filter functions} $f$, $g$ (Fig.~\ref{fig:filters}) send an eigenvalue to a flag state
$\ket{h(\lambda)}=\sqrt{1-f^2-g^2}\ket{\mathit{nothing}}+f\ket{\mathit{well}}+g\ket{\mathit{ill}}$, and the
paper's Theorem~1 bounds the error of the flagged output in terms of the total evolution time $t=Nt_0$.

\begin{theorem}[Theorem~1 of the HHL paper, with constants]\label{thm:thm1}
Let $A$ be Hermitian with $|t_0\lambda|\le\pi/2$ on its spectrum, let the flag gates prepare
$\ket{h(\tilde\lambda_k)}$ for each clock value~$k$, and let $c=\pi^2\sqrt{4+2\pi^2}$. For every $b$, the
flagged output is within $3c\,\kappa\norm b/t$ of the ideal one, and the states post-selected on
$\{\mathit{well},\mathit{ill}\}$ (or on $\mathit{well}$, for $b$ in the well-conditioned subspace) are within
$12c\,\kappa/t$ of the ideal ones.\isalink{Algorithms/HHL_Filter_Bounds.thy}{330}
\end{theorem}

\begin{theorem}[run time]\label{thm:runtime}
Let the spectrum of $A$ lie in $[1/\kappa,1]$, let $t_0\le\pi/2$, $\varepsilon\kappa\le1$ and
$Nt_0\ge12c\,\kappa/\varepsilon$, and let $\norm b=1$. Then the post-selected output is within
$\varepsilon$ of $\psi\otimes A^{-1}b/\norm{A^{-1}b}\otimes\ket{\mathit{well}}$, its amplitude is at least
$1/(4\kappa)$, and amplitude amplification with $k\le\pi\kappa$ rounds and $(2k+1)\cdot2(N-1)$ calls of the
evolution succeeds with probability at least $1-\norm y^2$, where $y$ is the post-selected part of
one run.\isalink{Algorithms/HHL_Runtime.thy}{302}
\end{theorem}

\noindent An oracle $V$ with $\norm{Vv-e^{iAt_0}v}\le\eta\norm v$ moves every output by at most the number
of calls times~$\eta$, so the theorems hold for a simulated evolution with an additional error
$2(N-1)\eta$. The rest of the paper is verified too: the reduction from circuits to matrix inversion
behind its optimality theorems (with $\kappa\le2T+1$ for a circuit of $T$ gates), rectangular and
preconditioned systems, and the estimation of $\langle x|M|x\rangle$ from samples. Several of its
statements turned out to be false (Section~\ref{sec:errata}).
